\section{Vigésimo segunda iteración de la Fase 5: Texto en Pantalla}
\label{sec:F5I22}

F5-I21 dejó fijada la regla «cada dato se edita solo en su pestaña natural». Esta iteración la completa para los botones y simplifica el formato: \textbf{el texto que se ve en el botón es un único dato, «Texto Pantalla» (\texttt{textoPantalla}), que se edita SOLO en el Setup y viaja desde allí a los cartuchos GWB/GWI y al simulador}. Desaparece el campo \texttt{alias} del GWS y el campo \texttt{id} de los cartuchos GWB/GWI de nueva escritura: dos nombres para el mismo dato confunden. Al cargar, los ficheros antiguos se siguen aceptando y se convierten al campo nuevo.

\subsection{Análisis del estado actual}
\label{sec:F5I22Analisis}

\begin{itemize}
	\item Tras F5-I21, los botones seguían editándose en dos sitios: la columna \emph{Alias} de la tabla del Setup y la tabla (Usado / Id / Tecla) de la pestaña Cartucho. El \texttt{id} del cartucho es simplemente el texto que se pinta en el botón (de hecho solo se usaba para los nombres de sistema A/B/R y para \emph{Start}), igual que el \texttt{alias} del GWS. Duplicar el concepto con dos nombres confunde.
	\item El botón \emph{Aplicar} del Setup sincronizaba además esos datos hacia la ficha del Cartucho: dos sitios editables para el mismo dato contradicen la regla de F5-I21.
	\item En la pestaña Cartucho, la tabla de botones era edición de un dato cuya pestaña natural es el Setup: quedaba fuera de sitio, como le pasó a la resolución antes de F5-I21.
\end{itemize}

\subsection{Diseño}
\label{sec:F5I22Diseno}

\begin{itemize}
	\item \textbf{Un solo campo, «Texto Pantalla» (\texttt{textoPantalla}).} El GWS declara cada botón como \texttt{\{numero, puerto?, tecla?, textoPantalla?\}} y los cartuchos GWB/GWI lo repiten como \texttt{\{numero, tecla?, textoPantalla?\}}. Sin \texttt{textoPantalla}, el botón muestra su número (o, por compatibilidad, el \texttt{alias} de un GWS antiguo).
	\item \textbf{Compatibilidad hacia atrás al cargar:} un GWS con \texttt{alias}, o un GWB/GWI con \texttt{id}, se leen y se convierten al campo nuevo; ya no se escriben nunca en los ficheros nuevos. Los ficheros de datos existentes del repositorio (\texttt{setup*.json}, los cartuchos de ejemplo y los de \texttt{pruebas/}) no se tocan: siguen con \texttt{alias}/\texttt{id} y funcionan por esta compatibilidad; se convertirán al guardarlos de nuevo.
	\item \textbf{La pestaña Cartucho pierde la tabla de botones.} La ficha no muestra ni edita botones: puerto, tecla y «Texto Pantalla» se editan en el Setup. Los botones que viajan en el GWB/GWI generado son una copia de los del Setup vigente, con sus nombres de campo y sus datos.
	\item \textbf{El botón \emph{Aplicar} del Setup deja de sincronizar botones:} la ficha del Cartucho ya no los recibe ni los guarda, así que no hay nada que sincronizar.
	\item \textbf{El texto «Start» sigue designando el arranque:} el botón que el cartucho llama «Start» (venga de \texttt{textoPantalla} o del \texttt{id} antiguo) es el que enciende el juego, sea cual sea su número físico.
	\item \textbf{La máquina derivada del cartucho (F5-I19) hereda el texto.} Al cargar un cartucho sin setup en memoria, la máquina derivada toma puerto, tecla y «Texto Pantalla» botón a botón, con precedencia «el cartucho manda»: lo que trae el cartucho (\texttt{textoPantalla}, o \texttt{id}/\texttt{alias} antiguos) manda sobre el setup previo, y lo que nadie declara se hereda del setup anterior.
	\item \textbf{El trío de sistema A/B/R:} la columna \emph{Botón} de la tabla del Setup ya nombra los botones 6, 7 y 8; el \texttt{textoPantalla} es opcional para ellos (vacío = su número).
\end{itemize}

\subsection{Plan de cambios}
\label{sec:F5I22Plan}

\begin{itemize}
	\item \texttt{baseSetup.js}: la columna \texttt{alias} pasa a ser «Texto Pantalla» (\texttt{colTextoPantalla}); al salvar, el JSON lleva \texttt{textoPantalla}; al cargar, se acepta como texto de partida el \texttt{textoPantalla} nuevo o el \texttt{alias} antiguo.
	\item \texttt{baseFicha.js}: desaparece la sección de UI de botones (la tabla Usado/Id/Tecla, los cursores y \texttt{\_pintarTablaBotones}); una función \texttt{\_botonesDelSetup()} aporta al GWB/GWI los botones del setup vigente. La ficha ya no valida ni guarda botones.
	\item \texttt{baseDirector.js}: \texttt{botonesDeSetup()} expone una copia de los botones del setup; \texttt{\_aplicarJSONSetup} y \texttt{\_cargarCartucho} dejan de pasar botones a la ficha (el \emph{Aplicar} deja de sincronizarlos); \texttt{\_aplicarBotonesCartucho} etiqueta cada botón físico con su \texttt{textoPantalla} (aceptando \texttt{alias}/\texttt{id} antiguos) y detecta «Start»; \texttt{\_setupDerivadoDe} propaga el \texttt{textoPantalla} con precedencia cartucho $>$ setup previo; \texttt{\_construirBotonFisico} usa \texttt{textoPantalla} y \texttt{\_normalizarNumeroBoton} resuelve números por texto (y por \texttt{alias} antiguo).
	\item \texttt{estilos.css}: el selector de la columna pasa de \texttt{colAlias} a \texttt{colTextoPantalla}.
	\item Sin cambios en el formato GWC (los botones no viajan en él) ni en el orden del tick.
\end{itemize}

\subsection{Pruebas}
\label{sec:F5I22Pruebas}

\begin{itemize}
	\item Al cargar un cartucho, la tabla de botones ya no aparece en la pestaña Cartucho (las cajas del grupo \texttt{director} lo verifican de paso); el botón «Start» sigue encendiendo el juego venga el texto del campo nuevo o del \texttt{id} antiguo (caja D3).
	\item Un GWS con \texttt{alias} y un GWB/GWI con \texttt{id} antiguos se cargan sin error y se convierten a \texttt{textoPantalla} (cajas D3 y T3).
	\item Las teclas que llegan al simulador vienen del GWS (etiqueta del botón físico, tecla efectiva); en la máquina derivada, la tecla del cartucho cargado sigue mandando sobre la del setup previo (caja T3 de F5-I17, M8).
	\item Los ficheros nuevos generados no contienen \texttt{alias} ni \texttt{id} en los botones (S2 y S5 lo verifican de paso al comparar salidas).
	\item S2 y S5 completas en verde; el cambio no toca el formato GWC ni el orden del tick.
\end{itemize}

\subsection{Criterios de aceptación}
\label{sec:F5I22Criterios}

\begin{enumerate}
	\item El campo de los botones se llama y se edita como «Texto Pantalla» (\texttt{textoPantalla}); no existe ya \texttt{alias} en el GWS ni \texttt{id} en GWB/GWI de nueva escritura.
	\item La pestaña Cartucho no edita (ni muestra) la tabla de botones; los botones se editan solo en el Setup (puerto, tecla y texto).
	\item El botón \emph{Aplicar} del Setup ya no sincroniza datos de botones.
	\item La información de botones se mantiene en los cartuchos guardados: el GWB/GWI llevan los botones del Setup con sus nombres de campo y sus datos, de modo que llegan al GWC y al simulador.
	\item «Start» sigue designando el botón de arranque, venga de \texttt{textoPantalla} o del \texttt{id} antiguo.
	\item Un GWS/GWB/GWI antiguo (con \texttt{alias}/\texttt{id}) se carga sin error y se convierte al campo nuevo.
	\item S2 y S5 en verde al cerrar la iteración.
\end{enumerate}
